% ═══════════════════════════════════════════════
%  MỞ ĐẦU
% ═══════════════════════════════════════════════

\section{MỞ ĐẦU}
\label{ch:intro}

% ───────────────────────────────────────────────
\subsection{Lý do chọn đề tài}
% ───────────────────────────────────────────────

\indent Trò chơi điện tử giữ một vị trí khác thường trong nghiên cứu trí tuệ nhân tạo:
chúng đồng thời là phòng thí nghiệm và là sản phẩm. Một môi trường game cho phép chạy
hàng triệu lần thử với chi phí gần bằng không và không có rủi ro nào, điều mà robot
thật hay hệ thống giao thông thật không bao giờ cho phép. Ở chiều ngược lại, chính
ngành công nghiệp game có nhu cầu thương mại trực tiếp với những nhân vật máy biết cư
xử thuyết phục. Hiếm lĩnh vực nào mà quãng đường từ một kết quả nghiên cứu tới một
tính năng bán được lại ngắn đến thế.

\indent Học tăng cường là hướng tiếp cận ăn khớp nhất với bối cảnh này. Nó không đòi
hỏi tập dữ liệu gán nhãn sẵn, thứ gần như không thể có cho câu hỏi ``nước đi nào là
đúng trong tình huống này'', mà để tác nhân tự rút ra chiến lược từ hậu quả của chính
những hành động nó đã thử.

\indent Khả năng của học tăng cường đã được chứng minh qua hàng loạt thành tựu mang
tính bước ngoặt. Một hệ thống học tăng cường có thể chơi thành thạo hàng chục trò chơi
điện tử cổ điển chỉ từ hình ảnh thô của màn hình \cite{mnih2015humanlevel}, có thể
đánh bại kỳ thủ cờ vây hàng đầu thế giới \cite{silver2016mastering}, hay phối hợp
nhiều tác nhân để chiến thắng trong những trò chơi chiến thuật phức tạp
\cite{openai2019dota, vinyals2019grandmaster}. Điểm chung của các hệ thống này là
không một hành vi nào trong số đó được lập trình sẵn: chúng đều nảy sinh từ quá trình
thử và sai của chính tác nhân.

\indent Tuy nhiên, một trong những rào cản lớn nhất khi nghiên cứu và ứng dụng học
tăng cường chính là việc xây dựng môi trường huấn luyện phù hợp. Các môi trường hai
chiều đơn giản hoặc các bộ dữ liệu chuẩn thường được dùng trong nghiên cứu học thuật
tuy thuận tiện nhưng lại không phản ánh được độ phức tạp của trò chơi thực tế. Một môi
trường ba chiều với mô phỏng vật lý chân thực, không gian quan sát phong phú và các
tình huống thay đổi liên tục mới thực sự thử thách khả năng học của tác nhân, đồng
thời gần với bài toán phát triển game thương mại hơn. Việc thiếu vắng những môi trường
như vậy, cùng với chi phí xây dựng chúng, khiến cho việc so sánh các thuật toán học
tăng cường trong bối cảnh game ba chiều vẫn còn là một khoảng trống đáng chú ý.

\indent Nền tảng Unity kết hợp với bộ công cụ ML-Agents cung cấp một hệ sinh thái đủ
mạnh để giải quyết vấn đề trên. Unity là một trong những công cụ phát triển game phổ
biến nhất hiện nay, hỗ trợ mô phỏng vật lý và dựng hình ba chiều ở mức chuyên nghiệp;
còn ML-Agents nối thẳng một cảnh Unity đang chạy với các thuật toán học tăng cường
được cài đặt trong hệ sinh thái Python.
Sự kết hợp này biến Unity thành một nền tảng lý tưởng để thiết kế các môi trường huấn
luyện đa dạng, từ bài toán đối kháng, hợp tác cho đến điều hướng tránh vật cản trong
không gian ba chiều.

\indent Xuất phát từ những lý do trên, tôi lựa chọn đề tài ``Học tăng cường
để tối ưu hóa khả năng tương tác của AI trong game'' với trọng tâm là xây dựng ba
môi trường game ba chiều mang tính chất khác nhau trên nền tảng Unity ML-Agents, sau
đó huấn luyện và so sánh một cách có hệ thống năm thuật toán tiêu biểu là PPO, SAC,
MA-POCA, MAPPO và DQN trên từng loại môi trường. Đích đến của đề tài không phải là tìm ra
thuật toán mạnh nhất - kết quả thực nghiệm cho thấy câu hỏi đó không có lời giải - mà
là chỉ ra được mỗi thuật toán thắng ở đâu và vì sao, đồng thời để lại một bộ môi trường
dùng lại được cho các nghiên cứu và dự án phát triển game về sau.

% ───────────────────────────────────────────────
\subsection{Mục tiêu nghiên cứu}

\indent Mục tiêu tổng quát của đề tài là nghiên cứu, xây dựng và đánh giá quy trình
ứng dụng học tăng cường để huấn luyện tác nhân trí tuệ nhân tạo trong môi trường game
ba chiều trên nền tảng Unity ML-Agents. Từ mục tiêu tổng quát đó, đề tài hướng đến năm
mục tiêu cụ thể.

\begin{itemize}
    \item \textbf{Xây dựng bộ môi trường thực nghiệm.} Đề tài xây dựng ba môi trường
    game ba chiều có tính chất khác biệt để đại diện cho ba dạng bài toán học tăng cường
    phổ biến: Football Table mô phỏng một bàn bi lắc ba chiều, đại diện cho bài toán đối
    kháng một đối một với không gian hành động liên tục; Capture The Flag đòi hỏi nhiều
    tác nhân phối hợp với nhau để cùng hoàn thành một mục tiêu chung; và Cross The Road
    buộc tác nhân điều hướng vượt qua dòng phương tiện di chuyển liên tục để đến đích.

    \item \textbf{Thiết kế quan sát, hành động và phần thưởng.} Đề tài thiết kế hệ thống
    quan sát, hành động và phần thưởng phù hợp với đặc thù của từng môi trường, sao cho
    tác nhân có đủ thông tin cần thiết để học được hành vi hiệu quả trong không gian ba
    chiều mà không rơi vào những hành vi bất thường ngoài mong muốn.

    \item \textbf{Huấn luyện và thu thập số liệu.} Đề tài triển khai huấn luyện năm thuật
    toán PPO, SAC, MA-POCA, MAPPO và DQN trên từng môi trường trong cùng điều kiện phần
    cứng và cùng ngân sách bước, tạo thành một lưới đầy đủ mười lăm lần chạy; kết quả
    được thu thập và phân tích theo bốn tiêu chí định lượng là mức năng lực đạt được,
    tốc độ hội tụ, độ ổn định của quá trình huấn luyện và độ dài lượt chơi.

    \item \textbf{So sánh và khuyến nghị lựa chọn thuật toán.} Trên cơ sở các kết quả thu
    được, đề tài so sánh và phân tích ưu nhược điểm của từng thuật toán ứng với từng loại
    bài toán, từ đó đưa ra khuyến nghị lựa chọn thuật toán phù hợp cho các cơ chế game
    khác nhau trong thực tiễn phát triển.

    \item \textbf{Xây dựng ứng dụng trình diễn.} Đề tài xây dựng một ứng dụng trình diễn
    tương tác hoàn chỉnh, trong đó các mô hình đã huấn luyện (định dạng ONNX) được triển
    khai trực tiếp vào game. Ứng dụng cung cấp menu lựa chọn môi trường, bốn chế độ chơi
    (người chơi tự điều khiển, tác nhân tự chơi, người chơi cùng tác nhân, tác nhân đấu
    mô hình ngôn ngữ lớn), cơ chế lựa chọn mô hình cho tác nhân và hai tab tra cứu số
    liệu huấn luyện được trích xuất trực tiếp từ nhật ký TensorBoard, qua đó khép kín quy
    trình từ huấn luyện đến trải nghiệm thực tế trong game.
\end{itemize}

% ───────────────────────────────────────────────
\subsection{Đối tượng và phạm vi nghiên cứu}

\indent Đối tượng nghiên cứu của đề tài gồm ba nhóm chính. Nhóm môi trường bao gồm ba
môi trường game ba chiều được xây dựng trên nền tảng Unity, gồm môi trường đối kháng
Football Table, môi trường Capture The Flag đa tác nhân và môi trường điều hướng Cross
The Road. Nhóm thuật toán bao gồm năm thuật toán học tăng cường được khảo sát, gồm PPO
(thuật toán tối ưu chính sách vùng lân cận), SAC (thuật toán actor-critic cực đại hóa
entropy), MA-POCA (thuật toán gán công lao cho đa tác nhân), MAPPO (mở rộng PPO cho
đa tác nhân với chia sẻ tham số) và DQN (thuật toán dựa trên giá trị với bộ đệm phát
lại). Nhóm công cụ bao gồm bản thân nền tảng Unity ML-Agents, cơ chế trainer mở rộng
dùng để bổ sung MAPPO và DQN, cùng giao diện lập trình cấp thấp bằng Python phục vụ
việc tích hợp thuật toán tùy chỉnh.

\indent Về phạm vi nội dung, đề tài tập trung vào thiết kế môi trường ba chiều, xây
dựng hệ thống quan sát, hành động, phần thưởng, huấn luyện tác nhân và đánh giá hiệu
năng thuật toán. Đề tài không đi sâu vào các khía cạnh tối ưu hóa đồ họa, âm thanh hay
hoàn thiện sản phẩm game thương mại. Ba môi trường được lựa chọn nhằm bao phủ ba dạng
bài toán học tăng cường điển hình là đối kháng một đối một, hợp tác đa tác nhân và
điều hướng đơn tác nhân.

\indent Về phạm vi thuật toán, năm thuật toán PPO, SAC, MA-POCA, MAPPO và DQN được chọn
dựa trên mức độ phổ biến trong nghiên cứu, tính sẵn có trong hệ sinh thái ML-Agents và
sự đa dạng về mặt kiến trúc, khi chúng đại diện cho năm hướng tiếp cận khác nhau lần
lượt là on-policy đơn tác nhân, actor-critic off-policy, đa tác nhân dựa trên attention,
đa tác nhân dựa trên chia sẻ tham số và phương pháp dựa trên giá trị. Vì DQN đòi hỏi
không gian hành động rời rạc, môi trường Football Table vốn dùng hành động liên tục
phải được dựng lại thành một bản riêng có hành động rời rạc dành cho lần chạy này.
Toàn bộ thực nghiệm được tiến hành trên một cấu hình máy
tính cố định để bảo đảm tính công bằng khi so sánh, chi tiết cấu hình này được trình
bày trong Mục~\ref{sub:hardware}.

% ───────────────────────────────────────────────
\subsection{Tổng quan tài liệu và khoảng trống nghiên cứu}

\indent Nghiên cứu về học tăng cường trong game đã trải qua nhiều giai đoạn phát triển.
Mnih và cộng sự \cite{mnih2015humanlevel} đề xuất DQN và lần đầu chứng minh rằng học
tăng cường sâu có thể đạt trình độ con người trong các trò chơi điện tử cổ điển chỉ từ
dữ liệu điểm ảnh thô, mở ra hướng ứng dụng học tăng cường trực tiếp vào môi trường
game; đây cũng là lý do DQN được giữ lại trong đề tài như một mốc so sánh lịch sử cho
nhóm phương pháp dựa trên giá trị. Tiếp nối
thành tựu đó, Silver và cộng sự \cite{silver2016mastering, silver2017mastering} đạt
bước đột phá với các hệ thống chơi cờ vây ở trình độ kiện tướng thế giới, khi kết hợp
học tăng cường với mạng nơ-ron sâu và tìm kiếm cây.

\indent Ở quy mô lớn hơn với môi trường phức tạp và nhiều tác nhân, các hệ thống chơi
trò chơi chiến thuật thời gian thực \cite{openai2019dota, vinyals2019grandmaster} đã
cho thấy học tăng cường có thể xử lý những bài toán với không gian trạng thái khổng
lồ, nhiều tác nhân hoạt động đồng thời và đòi hỏi chiến lược dài hạn. Về phía thuật
toán, Schulman và cộng sự \cite{schulman2017proximal} đề xuất PPO, một phương pháp
on-policy chấp nhận đánh đổi hiệu quả mẫu để lấy sự ổn định và tính dễ triển khai,
trong khi Haarnoja và cộng sự
\cite{haarnoja2018soft} giới thiệu SAC với cơ chế điều hòa entropy giúp cân bằng giữa
khám phá và khai thác, đặc biệt hiệu quả với không gian hành động liên tục. Trong lĩnh
vực đa tác nhân, Cohen và cộng sự \cite{cohen2022mapoca} phát triển MA-POCA với cơ chế
phê bình tập trung nhằm giải quyết bài toán gán công lao trong môi trường hợp tác,
trong khi Yu và cộng sự \cite{yu2022surprising} chứng minh rằng MAPPO - một mở rộng đơn
giản của PPO sang nhiều tác nhân với chia sẻ tham số - đạt hiệu năng cạnh tranh hoặc
vượt trội trên nhiều bộ benchmark hợp tác chuẩn.

\indent Về mặt công cụ, Juliani và cộng sự \cite{juliani2020unity} giới thiệu Unity
ML-Agents như một hệ sinh thái mã nguồn mở tích hợp sẵn khả năng xây dựng môi trường
ba chiều với mô phỏng vật lý chân thực, hỗ trợ huấn luyện tác nhân trực tiếp và được
sử dụng rộng rãi trong cả nghiên cứu lẫn phát triển game.

\indent Điều khiến các kết quả trên khó dùng lại nằm ở chỗ chúng gần như luôn được đo
trên những bộ benchmark đã chuẩn hóa: lưới hai chiều, vector trạng thái rút gọn, động
lực học đơn giản. Một thứ hạng thuật toán rút ra từ đó nói được rất ít về hành vi của
cùng thuật toán ấy khi phải đương đầu với va chạm vật lý liên tục, quan sát một phần
qua cảm biến tia và phần thưởng thưa của một cảnh ba chiều thật. Mỗi công trình lại
thường chỉ khảo sát một hoặc hai thuật toán trên bộ môi trường riêng của mình, nên
việc đặt PPO, SAC, MA-POCA, MAPPO và DQN lên cùng một bàn cân gần như không thực hiện được
chỉ bằng cách đối chiếu tài liệu đã công bố. Đề tài này chấp nhận tự dựng lấy bàn cân
đó: ba môi trường ba chiều thuộc ba dạng bài toán khác nhau, chung một cấu hình phần
cứng và chung một quy trình đo.

% ───────────────────────────────────────────────
\subsection{Phương pháp nghiên cứu}
\indent Đề tài sử dụng kết hợp ba nhóm phương pháp nghiên cứu bổ trợ lẫn nhau. Phương
pháp nghiên cứu lý thuyết được dùng để tổng hợp và phân tích các tài liệu khoa học về
học tăng cường, học tăng cường sâu và năm thuật toán PPO, SAC, MA-POCA, MAPPO và DQN,
dựa trên các bài báo gốc cùng tài liệu kỹ thuật của Unity ML-Agents; đây là cơ sở cho
việc thiết kế môi trường và lựa chọn tham số huấn luyện.

\indent Phương pháp thực nghiệm được sử dụng để xây dựng ba môi trường game ba chiều
trong Unity với đầy đủ hệ thống vật lý, quan sát, hành động và phần thưởng, sau đó
triển khai và huấn luyện từng thuật toán trên mỗi môi trường trong cùng điều kiện phần
cứng và số bước huấn luyện. Trong quá trình này, dữ liệu huấn luyện được thu thập và
trực quan hóa theo thời gian thực bằng công cụ TensorBoard.

\indent Cuối cùng, phương pháp so sánh và đánh giá được dùng để đối chiếu hiệu năng
của các thuật toán dựa trên bốn tiêu chí định lượng là mức năng lực đạt được, tốc độ
hội tụ, độ ổn định huấn luyện và độ dài lượt chơi, kết hợp với phân tích
định tính về mức độ phù hợp của từng thuật toán với từng dạng môi trường.

% ───────────────────────────────────────────────
\subsection{Ý nghĩa khoa học và thực tiễn}

\indent Về mặt khoa học, đề tài đề xuất và hiện thực hóa một quy trình xây dựng môi
trường game ba chiều phục vụ huấn luyện học tăng cường trên nền tảng Unity, có thể
được tái sử dụng như một bộ môi trường tham chiếu mở cho các nghiên cứu tiếp theo.
Bên cạnh đó, bộ thực nghiệm so sánh có hệ thống giữa PPO, SAC, MA-POCA, MAPPO và DQN
trên ba dạng môi trường đặc trưng góp phần làm rõ điều kiện phù hợp của từng thuật toán
trong bối
cảnh game ba chiều. Ngoài ra, quy trình bổ sung thuật toán mới vào ML-Agents dưới dạng
trainer mở rộng, cùng cách kết nối một thuật toán off-policy với Unity thông
qua giao diện lập trình cấp thấp bằng Python, cũng được trình bày chi tiết, mở ra một
hướng đi có thể tái sử dụng cho các nghiên cứu cần thuật toán nằm ngoài bộ tích hợp
sẵn của Unity.

\indent Về mặt thực tiễn, ba môi trường ba chiều được xây dựng trong đề tài có thể
được sử dụng trực tiếp làm nền tảng thử nghiệm trí tuệ nhân tạo trong các dự án game
Unity khác. Kết quả so sánh thuật toán giúp các nhà phát triển có cơ sở lựa chọn thuật
toán học tăng cường phù hợp với từng cơ chế game, chẳng hạn cơ chế đối kháng, hợp tác
hay điều hướng, mà không phải tiến hành lại toàn bộ thực nghiệm từ đầu. Quy trình
thiết kế hàm phần thưởng và không gian quan sát trong môi trường ba chiều được trình
bày trong đề tài cũng có giá trị tham khảo cho các nhóm phát triển game có nhu cầu
tích hợp trí tuệ nhân tạo.

% ───────────────────────────────────────────────
\subsection{Bố cục đồ án}
\indent Ngoài phần Mở đầu và Kết luận, nội dung chính của đồ án được tổ chức thành bốn
chương. Chương~\ref{ch:theory} trình bày cơ sở lý thuyết, bao gồm tổng quan về trí tuệ
nhân tạo trong game, nền tảng lý thuyết của học tăng cường với quá trình quyết định
Markov, chính sách và hàm giá trị, các kiến trúc học tăng cường sâu, cùng đặc điểm của
năm thuật toán PPO, SAC, MA-POCA, MAPPO và DQN; chương này cũng giới thiệu nền tảng Unity
ML-Agents và các giới hạn tích hợp liên quan đến việc triển khai thuật toán tùy chỉnh.

\indent Chương~\ref{ch:design} trình bày chi tiết quá trình thiết kế ba môi trường
game ba chiều trong Unity, bao gồm không gian quan sát, không gian hành động và hàm
phần thưởng của từng môi trường; đồng thời mô tả phương pháp đưa các thuật toán nằm
ngoài bộ huấn luyện tích hợp sẵn của ML-Agents vào thực nghiệm, thiết lập thực nghiệm
và các tiêu chí đánh giá được sử dụng. Chương~\ref{ch:results} trình bày kết quả huấn
luyện của từng thuật toán trên từng môi trường thông qua biểu đồ và bảng số liệu, tiến
hành so sánh tổng hợp, phân tích ưu nhược điểm của từng thuật toán theo từng dạng bài
toán và thảo luận về khả năng ứng dụng thực tiễn cùng hướng phát triển tiếp theo.

\indent Cuối cùng, Chương~\ref{ch:app} trình bày quá trình xây dựng ứng dụng trình
diễn tương tác: kiến trúc hệ thống lựa chọn chế độ chơi, giao diện menu, cơ chế nạp
mô hình ONNX tại thời gian chạy, đường ống đưa số liệu huấn luyện từ TensorBoard vào
giao diện trong game, cách ánh xạ bốn chế độ chơi (người chơi, tác nhân tự chơi, người
chơi cùng tác nhân, tác nhân đấu mô hình ngôn ngữ lớn) vào từng môi trường, cùng các
vấn đề kỹ thuật phát sinh và kết quả kiểm thử ứng dụng.
